\documentclass[12pt]{article}
\usepackage[utf8]{inputenc}
\usepackage[T1]{fontenc}
\usepackage{hyperref}
\title{Kwandu Language Resources for K-12 Curriculum Development}
\author{Compiled by Claude Sonnet 5.5 (claude-sonnet-5-5) for Andr\'as Kornai}
\date{}
\begin{document}
\maketitle

\section{Introduction}
Kwandu (ISO 639-3 xdo; Glottocode kwan1289) is a Bantu language of southwestern Angola. Joshua Project describes it as a Kavango-Southwest Bantu language, also called Mashi, with North Kwandu and South Kwandu dialects. It places the homeland in semi-arid mountain areas around Serra da Neve, with places near Kamukuio and Mamue. Do not confuse Kwandu with Kuvale, which is a separate language.

\section{Primary Educational Resources}
I found no primers, textbooks, orthography or school programmes for Kwandu.

\section{Linguistic and Reference Materials}
The searches turned up no grammar or dictionary of Kwandu. Related material on the neighbouring Kuvale language exists, but it concerns a different language:
\begin{itemize}
\item Wikipedia, ``Kuvale language'', for context only. \url{https://en.wikipedia.org/wiki/Kuvale_language}
\end{itemize}

\section{Religious and Community Resources}
\begin{itemize}
\item Joshua Project, ``Kwandu in Angola'', a people-group profile. \url{https://joshuaproject.net/people_groups/21064/AO}
\end{itemize}
I saw no scripture translation URL.

\section{Digital and Audio Resources}
None found.

\section{Curriculum Development Context}
Very little is documented online for Kwandu. A curriculum would need fieldwork with the community, starting with an orthography and word lists. Angolan schooling is in Portuguese. I found no evidence of mother-tongue instruction in Kwandu.
\end{document}